\documentclass[11pt,letterpaper]{article}

\usepackage[margin=1in]{geometry}
\usepackage{graphicx}
\usepackage{booktabs}
\usepackage{amsmath}
\usepackage{hyperref}
\usepackage{xcolor}
\usepackage{caption}
\usepackage{float}

\hypersetup{colorlinks=true, linkcolor=blue, citecolor=blue, urlcolor=blue}

% --- Local macros ---
\newcommand{\pp}{\mbox{$p$+$p$}}
\newcommand{\sqs}{\mbox{$\sqrt{s}$}}
\newcommand{\pT}{\ensuremath{p_{\text{T}}}}
\newcommand{\ET}{\ensuremath{E_{\text{T}}}}
\newcommand{\zvtx}{\ensuremath{z_{\mathrm{vtx}}}}
\newcommand{\abseta}{\ensuremath{|\eta|}}
\newcommand{\Rcemc}{\ensuremath{R_{\text{CEMC}}}}
\newcommand{\vtxeff}{\ensuremath{\varepsilon_{\text{vtx+MBD}}}}
\newcommand{\fcorr}{\ensuremath{F_{\text{corr}}}}
\newcommand{\fdouble}{\ensuremath{f_{\text{double}}}}

\graphicspath{{../../plotting/figures/pileup_vtx/}}

\title{%
  \vspace{-1cm}
  {\large \textbf{sPHENIX Internal Note}}\\[0.5cm]
  {\LARGE \textbf{Pileup Correction to Vertex + MBD Acceptance\\
  (Luminosity Note Section~3)}}\\[0.3cm]
  {\large Isolated Photon Cross-Section in \pp{} at $\sqs = 200$~GeV}
}
\author{sPHENIX Collaboration}
\date{\today}

\begin{document}

\maketitle

\begin{abstract}
This note quantifies the pileup correction to the vertex + MBD
coincidence acceptance in the sPHENIX isolated photon cross-section
measurement (PPG12). In double-interaction events, the overlay
collision provides additional MBD hits and shifts the reconstructed
vertex toward the beam center, increasing the effective acceptance
from $\vtxeff = 0.68$ (single interaction) to $0.81$ (double
interaction) --- a 19\% relative increase. After Poisson weighting
with the measured double-interaction fractions, the net correction
to the cross-section is $-2.1\%$ for 0~mrad ($\fdouble = 11.1\%$)
and $-0.7\%$ for 1.5~mrad ($\fdouble = 3.9\%$), with a mild
$\pT$-dependence rising from $+2\%$ at 8~GeV to $+3.5\%$ at
30~GeV. This correction is currently not applied in the production
efficiency chain, which uses single-interaction MC only.
\end{abstract}

\tableofcontents
\newpage

%% ============================================================
\section{Motivation}
\label{sec:motivation}

The sPHENIX luminosity note (Section~3) describes a correction
needed when the analysis requires both an MBD coincidence and a
z-vertex selection. In PPG12, events must satisfy:
\begin{itemize}
  \item MBD N\&S $\geq 1$ (at least one hit on each side)
  \item $|z_{\text{vtx}}^{\text{reco}}| < 60$~cm
\end{itemize}

The production efficiency chain (\texttt{RecoEffCalculator\_TTreeReader.C}
$\to$ \texttt{CalculatePhotonYield.C}) computes \vtxeff{} from
single-interaction MC:
\begin{equation}
  \vtxeff = \frac{N(\text{truth photons passing MBD + vertex cut})}
                 {N(\text{truth photons with } |z_{\text{vtx}}^{\text{truth}}| < 60~\text{cm})}
  \label{eq:vtxeff}
\end{equation}

In data, pileup events modify this acceptance in two ways:
\begin{enumerate}
  \item \textbf{MBD rescue:} If the signal collision fires only one
    MBD side (or neither), the overlay collision can provide the
    missing hits, satisfying the coincidence requirement.
  \item \textbf{Vertex centering:} The MBD-reconstructed vertex is
    approximately the average of two collision vertices. This pulls
    the vertex toward zero, increasing the fraction passing the
    $|z_{\text{vtx}}| < 60$~cm cut.
\end{enumerate}

The correction is:
\begin{equation}
  \vtxeff^{\text{blended}} = (1 - \fdouble) \times \text{Pass}_1
    + \fdouble \times \text{Pass}_2
  \label{eq:blend}
\end{equation}
where Pass$_1$ is the single-interaction acceptance and Pass$_2$ is
the double-interaction acceptance, and $\fdouble$ is the event-level
double-interaction fraction from Poisson statistics.

%% ============================================================
\section{Method}
\label{sec:method}

\subsection{Samples}
\label{sec:samples}

Single-interaction MC uses the three standard photon-jet samples,
weighted by production cross-section:
\begin{itemize}
  \item \texttt{photon5} (0--14~GeV), weight $= \sigma_5 / \sigma_{20} = 1122$
  \item \texttt{photon10} (14--30~GeV), weight $= \sigma_{10} / \sigma_{20} = 53.2$
  \item \texttt{photon20} (30+~GeV), weight $= 1.0$ (reference)
\end{itemize}

Double-interaction MC uses the full GEANT4 double-interaction sample:
\begin{itemize}
  \item \texttt{photon10\_double} --- two independent pp collisions
    simulated per event, with the photon from the primary collision
    in the 10--100~GeV truth $\pT$ range.
\end{itemize}

Two million events are processed per sample.

\subsection{Truth photon selection}
\label{sec:selection}

Standard PPG12 truth criteria:
$\texttt{pid} = 22$, $\texttt{photonclass} < 3$,
$E_\text{T}^\text{iso,truth} < 4$~GeV (cone $R = 0.3$),
$\abseta < 0.7$, $8 \leq \pT < 36$~GeV.

The denominator of Eq.~(\ref{eq:vtxeff}) requires
$|z_{\text{vtx}}^{\text{truth}}| < 60$~cm (same as
\texttt{RecoEffCalculator\_TTreeReader.C:1694}).

\subsection{Acceptance categories}
\label{sec:categories}

For each truth photon in the denominator, we check four acceptance
conditions independently:
\begin{enumerate}
  \item \textbf{MBD + vertex (combined):} MBD N\&S $\geq 1$ AND
    $|z_{\text{vtx}}^{\text{reco}}| < 60$~cm --- this is \vtxeff{}
    used in the cross-section formula.
  \item \textbf{MBD only:} MBD N\&S $\geq 1$ (no reco vertex cut).
  \item \textbf{Vertex only:} $|z_{\text{vtx}}^{\text{reco}}| < 60$~cm
    (no MBD requirement).
  \item \textbf{MBD hit decomposition:} classified into both N\&S,
    N only, S only, neither.
\end{enumerate}

\subsection{Pileup fractions}
\label{sec:fractions}

Event-level double-interaction fractions from
\texttt{calc\_pileup\_range.C} (Poisson statistics, lumi-weighted):
\begin{itemize}
  \item 0~mrad: $\fdouble = 11.1\%$ (event-level; 22.4\%
    cluster-weighted)
  \item 1.5~mrad: $\fdouble = 3.9\%$ (event-level; 7.9\%
    cluster-weighted)
\end{itemize}

Triple+ interactions are folded into the double fraction
(Prob$_3 < 0.3\%$ for $\mu < 0.3$).

%% ============================================================
\section{Results}
\label{sec:results}

\subsection{Single vs.\ double-interaction acceptance}
\label{sec:comparison}

Figure~\ref{fig:vtxeff} shows the per-$\pT$ vertex+MBD acceptance
for single-interaction (blue) and double-interaction (red) MC.
The double-interaction acceptance is systematically higher across
all $\pT$ bins, with a ratio of $\sim$1.19 (Figure~\ref{fig:ratio}).

\begin{figure}[htbp]
\centering
\includegraphics[width=0.85\textwidth]{vertex_eff_single_vs_double.pdf}
\caption{Per-$\pT$ vertex+MBD acceptance for single-interaction
  (blue) and double-interaction (red) photon MC. Error bars are
  binomial. The double-interaction acceptance is 19\% higher on
  average due to pileup MBD rescue and vertex centering.}
\label{fig:vtxeff}
\end{figure}

\begin{figure}[htbp]
\centering
\includegraphics[width=0.85\textwidth]{vertex_eff_ratio.pdf}
\caption{Top: vertex+MBD acceptance for both interaction types.
  Bottom: ratio (double/single), showing a $\sim$1.19 enhancement
  that increases slightly with $\pT$ (from 1.17 at 8~GeV to
  1.28 at 30~GeV).}
\label{fig:ratio}
\end{figure}

\subsection{Acceptance breakdown}
\label{sec:breakdown}

Figure~\ref{fig:breakdown} separates the MBD and vertex contributions:
\begin{itemize}
  \item \textbf{MBD acceptance:} increases from 69\% (single) to
    92\% (double) --- a 33\% relative increase. The overlay
    collision rescues events where the signal fired only one MBD
    side or neither.
  \item \textbf{Vertex acceptance:} increases from 75\% (single) to
    83\% (double) --- an 11\% relative increase. The averaged vertex
    is narrower than individual vertices.
\end{itemize}

The MBD rescue effect dominates the correction.

\begin{figure}[htbp]
\centering
\includegraphics[width=\textwidth]{acceptance_breakdown.pdf}
\caption{Acceptance breakdown into MBD-only (triangles up),
  vertex-only (triangles down), and combined (squares) for
  single (left) and double (right) interaction MC.}
\label{fig:breakdown}
\end{figure}

\subsection{MBD hit decomposition}
\label{sec:mbd_hits}

Figure~\ref{fig:mbd_decomp} shows the stacked MBD hit fractions.
In single-interaction MC, $\sim$8\% of truth photons produce no MBD
hits (``neither''), consistent with the GEANT4 prediction
($P_\text{neither} = 8.8\%$). In double-interaction MC, this drops
to $< 1\%$ because the overlay almost always provides some MBD
activity.

\begin{figure}[htbp]
\centering
\includegraphics[width=\textwidth]{mbd_hit_decomposition.pdf}
\caption{MBD hit decomposition for truth photons within the vertex
  acceptance. ``Both N\&S'' = MBD coincidence satisfied;
  ``N only'' / ``S only'' = single-arm fire; ``Neither'' = no
  MBD hits. The double-interaction overlay nearly eliminates the
  ``neither'' category.}
\label{fig:mbd_decomp}
\end{figure}

\subsection{Correction factors}
\label{sec:correction}

Figure~\ref{fig:correction} shows the per-$\pT$ correction factor
$\vtxeff^{\text{blended}} / \vtxeff^{\text{single}}$ for both
crossing-angle periods.

\begin{figure}[htbp]
\centering
\includegraphics[width=0.85\textwidth]{correction_factor_vs_pt.pdf}
\caption{Per-$\pT$ correction factor for 0~mrad (red, $\fdouble =
  11.1\%$) and 1.5~mrad (blue, $\fdouble = 3.9\%$). The correction
  increases with $\pT$ because the single-interaction acceptance
  drops faster at high $\pT$ (forward photons more sensitive to
  vertex shifts), while double-interaction rescues proportionally
  more events.}
\label{fig:correction}
\end{figure}

%% ============================================================
\section{Impact on the Cross-Section}
\label{sec:impact}

The correction increases \vtxeff{} in the efficiency denominator,
which \emph{decreases} the measured cross-section:
\begin{equation}
  \frac{d\sigma}{d\pT}\bigg|_{\text{corrected}} =
    \frac{d\sigma}{d\pT}\bigg|_{\text{current}}
    \times \frac{\vtxeff^{\text{single}}}{\vtxeff^{\text{blended}}}
\end{equation}

\begin{table}[htbp]
\centering
\caption{Integrated correction factors and cross-section impact.}
\label{tab:impact}
\begin{tabular}{lccc}
\toprule
\textbf{Period} & $\fdouble$ & \textbf{Integrated correction}
  & \textbf{Cross-section change} \\
\midrule
0 mrad    & 11.1\% & $+2.1\%$ & $-2.1\%$ \\
1.5 mrad  & 3.9\%  & $+0.7\%$ & $-0.7\%$ \\
\bottomrule
\end{tabular}
\end{table}

The $\pT$-dependent correction ranges from $+2.0\%$ at $8$~GeV to
$+3.7\%$ at $26$~GeV for the 0~mrad period.

%% ============================================================
\section{Discussion}
\label{sec:discussion}

\subsection{Comparison with other systematics}

The Section~3 correction is smaller than the dominant PPG12
systematics:
\begin{itemize}
  \item Purity systematic: 5--10\%
  \item Energy scale: 3--5\%
  \item $\sigma_\text{MBD}$ uncertainty: $\sim$7\% (flat)
\end{itemize}

For the 1.5~mrad nominal result, the $+0.7\%$ correction is
negligible. For 0~mrad or combined, the $+2.1\%$ (rising to
3.5\% at high $\pT$) is comparable to the energy scale
systematic and should be applied or assigned.

\subsection{Limitations}

\begin{enumerate}
  \item Only \texttt{photon10\_double} exists as full GEANT
    double-interaction MC. The correction uses the same
    Pass$_2$ for all $\pT$, though the actual correction should
    be computed per truth $\pT$ window (photon5/10/20).
  \item Triple+ interactions ($n \geq 3$) are folded into the
    double fraction. At $\mu \approx 0.12$ (0~mrad),
    Prob$_3 \approx 0.3\%$ --- negligible.
  \item The correction is computed as a single average over all
    runs, not per-segment. The $\mu$ variation within a fill
    introduces a $< 0.5\%$ residual uncertainty (Jensen's inequality).
  \item The double-interaction vertex distribution is from GEANT4
    simulation; any data-MC discrepancy in the beam vertex profile
    propagates.
\end{enumerate}

%% ============================================================
\section{Recommendation}
\label{sec:recommendation}

\begin{enumerate}
  \item \textbf{1.5~mrad (nominal):} No correction needed.
    Assign $+0.7\%$ as a one-sided systematic if desired.
  \item \textbf{0~mrad or combined:} Apply the per-$\pT$
    correction from Table in Figure~\ref{fig:correction},
    or assign $+2.1\%$ (integrated) as a systematic.
    The $\pT$-dependent correction table is available in
    \texttt{results/pileup\_vertex\_acceptance.csv}.
  \item \textbf{Implementation:} Multiply the current
    \texttt{vertex\_eff} in \texttt{CalculatePhotonYield.C:1009}
    by the per-$\pT$ correction factor for the appropriate
    crossing-angle period.
\end{enumerate}

%% ============================================================
\section*{Files}
\begin{tabular}{ll}
\toprule
\textbf{File} & \textbf{Purpose} \\
\midrule
\texttt{efficiencytool/study\_pileup\_vertex\_acceptance.py} & Analysis script \\
\texttt{plotting/plot\_pileup\_vertex\_correction.py} & Plotting script \\
\texttt{results/pileup\_vertex\_acceptance.root} & ROOT output \\
\texttt{results/pileup\_vertex\_acceptance.csv} & Per-$\pT$ table \\
\texttt{plotting/figures/pileup\_vtx/*.pdf} & Figures \\
\bottomrule
\end{tabular}

\end{document}
